\chapter{Beyond Softmax Attention: Linear Attention, State-Space Models, Hybrids, and Sparse Experts}\label{ch:attnalt}

Softmax attention has two costs that grow with context. Training does $O(T^2)$ work per layer; FlashAttention
(\cref{sec:attention:flash}) removes the $O(T^2)$ \emph{memory}, but not the FLOPs. Decoding reads a KV cache
that grows linearly with every token generated, so long-context inference is bounded by memory capacity and
bandwidth (\cref{ch:inference}). For a decade of chat-length contexts these costs were tolerable.
Reasoning traces of tens of thousands of tokens, agents with long tool histories, and million-token
documents have changed that, and several recent production models now replace most of their
attention layers with something cheaper.

This chapter covers that ``something cheaper''. We start from the incremental toolkit: windows,
sparsity and compressed caches. We then derive linear attention and show that it has three equivalent
computational forms, add decay and gating to reach the modern gated linear-attention and Mamba-2
family, add the delta rule, and connect all of it to state-space models and their parallel scans. We
then see how hybrids combine the pieces. A final section revisits mixture-of-experts layers with the
design details that \cref{sec:modern:moe} left out. Every operator comes with code in
\code{llmsys/linear\_attn.py} and \code{llmsys/ssm.py} that is verified to be equivalent across its forms,
because for these models the \emph{form} of the computation is the systems story.

\begin{figure}[t]
\centering
\resizebox{\linewidth}{!}{%
\begin{tikzpicture}[
  st/.style={draw=sysblue, fill=sysbluebg, rounded corners=2pt, minimum width=2.25cm,
    minimum height=1.05cm, align=center, font=\small, text width=2.15cm},
  num/.style={circle, fill=sysblue, text=white, font=\scriptsize\bfseries, inner sep=1.2pt},
  arr/.style={-{Stealth[length=2mm]}, thick, draw=black!60}]
\node[st] (a) at (0,0) {Why go\\beyond softmax};
\node[st] (b) at (2.8,0) {Windows \&\\sparse attention};
\node[st] (c) at (5.6,0) {Linear attention:\\three forms};
\node[st] (d) at (8.4,0) {Decay \&\\gating};
\node[st] (e) at (11.2,0) {Delta rule};
\node[st, fill=keygreenbg, draw=keygreen] (f) at (11.2,-1.8) {State-space\\models \& scans};
\node[st, fill=keygreenbg, draw=keygreen] (g) at (8.4,-1.8) {Hybrids};
\node[st, fill=keygreenbg, draw=keygreen] (h) at (5.6,-1.8) {MLA \& compressed caches};
\node[st, fill=fmemory!20, draw=napkin] (i) at (2.8,-1.8) {Sparse experts revisited};
\node[st, fill=labpurplebg, draw=labpurple] (j) at (0,-1.8) {Lab 8};
\foreach \n/\k in {a/1,b/2,c/3,d/4,e/5,f/6,g/7,h/8,i/9,j/10}
  \node[num] at (\n.north west) {\k};
\draw[arr] (a)--(b); \draw[arr] (b)--(c); \draw[arr] (c)--(d); \draw[arr] (d)--(e);
\draw[arr] (e)--(f); \draw[arr] (f)--(g); \draw[arr] (g)--(h); \draw[arr] (h)--(i); \draw[arr] (i)--(j);
\end{tikzpicture}}
\caption{Roadmap of this chapter: from cheap restrictions of softmax attention (top left), through
recurrent sequence mixers, to the hybrid and sparse architectures built from them.}
\label{fig:attnalt:roadmap}
\end{figure}

% =====================================================================
\section{Why look beyond softmax attention?}\label{sec:attnalt:why}
\covers{CS336 L4}

\begin{napkinmath}[The KV cache at long context]
Qwen3-8B has $L=36$ layers and 8 KV heads of dimension 128. Its bf16 KV cache costs
$2\cdot36\cdot8\cdot128\cdot2\,\text{B}=147{,}456$\,B per token. At a 128k-token context
($131{,}072$ tokens) a \emph{single} sequence needs $147{,}456\times131{,}072\approx19.3$\,GB, more than the
16\,GB of weights. Every decoded token must read all of it, which takes about 5.8\,ms at 3.35\,TB/s, so
attention alone caps one sequence at under 175 tokens/s before the weights are even counted. Batching does
not help, since each sequence brings its own cache. This is the pressure behind every technique in this
chapter.
\end{napkinmath}

There are three broad responses:
\begin{enumerate}
\item \textbf{Shrink what is cached} without changing what attention computes: GQA/MQA
  (\cref{sec:modern:gqa}), multi-head latent attention (\cref{sec:attnalt:mla}), and KV quantisation
  (\cref{ch:efficient}).
\item \textbf{Attend to fewer keys}: sliding windows, sinks, block-sparse patterns and learned top-$k$
  selection (\cref{sec:attnalt:sparse}). The cache can then be bounded or pruned.
\item \textbf{Replace the operator} with one whose state does not grow with $T$: linear attention, its
  gated and delta-rule descendants, and state-space models (\cref{sec:attnalt:linear,%
  sec:attnalt:gating,sec:attnalt:delta,sec:attnalt:ssm}). These usually appear in \emph{hybrids} that keep
  a few softmax layers (\cref{sec:attnalt:hybrid}).
\end{enumerate}

% =====================================================================
\section{Windows, sinks, and sparse attention}\label{sec:attnalt:sparse}
\covers{CS336 L3; CS336 L4}

\paragraph{Sliding windows.} A sliding-window layer lets query $i$ see only keys $j\in(i-w,i]$, which makes
compute $O(Tw)$ and lets the cache be a ring buffer of $w$ entries. Information still travels further than
$w$ because each layer extends the receptive field by $w$, but precise long-range retrieval suffers. Hence
the now-standard \emph{interleaving}: most layers are local and every $k$-th layer is global. Gemma 2
alternates local (window 4096) and global layers; Gemma 3 uses five local layers per global one; Cohere
reported a 3:1 local:global pattern in which the global layers use no positional encoding at all, so that
RoPE handles short-range order and NoPE layers handle long-range retrieval.

\paragraph{Attention sinks.} When a window slides past the first tokens of a sequence, quality can collapse.
Xiao et al.~(2023) traced this to the model using the first few tokens as a place to park attention mass:
softmax weights must sum to one, so heads with nothing useful to attend to dump their weight on early
tokens. Keeping those few \emph{sink} tokens permanently in the cache, alongside the window, restores
stable streaming generation.
The cache then has a fixed layout of $s$ sink slots followed by a ring buffer of $w$ recent slots, and the
index arithmetic is all a serving engine needs:

\begin{lstlisting}[caption={Slot assignment for a sink-plus-window cache of fixed size $s+w$.},label={lst:attnalt:sink}]
def cache_slot(pos, n_sink=4, window=4092):
    # where token `pos` lives in a cache with n_sink fixed slots + a ring of `window`
    if pos < n_sink:
        return pos                                    # sinks are never evicted
    return n_sink + (pos - n_sink) % window           # ring buffer overwrites the oldest

def live_positions(cur, n_sink=4, window=4092):
    # token positions visible to the query at position `cur`
    start = max(n_sink, cur - window + 1)
    return list(range(min(n_sink, cur + 1))) + list(range(start, cur + 1))
\end{lstlisting}
StreamingLLM assigns RoPE positions by slot order within the cache rather than by original text position,
so relative distances stay inside the trained range. That is a detail an engine must get right when it
applies rotations at cache-write time.


\paragraph{Learned sparsity.} Block-sparse and strided patterns (Child et al., 2019; Longformer's global
tokens) fix the pattern in advance. A newer line \emph{selects} keys per query: a cheap scorer (an
``indexer'') ranks all past tokens, and full attention runs only over the top-$k$. DeepSeek reported such a
scheme (DeepSeek Sparse Attention, in V3.2), applied by continued training of a dense model, which turns the
attention cost into $O(Tk)$ plus a much cheaper $O(T^2)$ indexing pass. The full cache is still stored, but
only $k$ entries are read per token.
\begin{napkinmath}[What does a top-$k$ indexer buy?]
Suppose an indexer with 4 heads of dimension 64, computed in FP8, scores every past token, and full attention
then runs over $k=2048$ selected tokens of a 128-head MLA-style layer. At $T=131{,}072$, the indexer does
about $T\cdot4\cdot64=3.4\times10^7$ multiply-adds per query token and reads $T\cdot64$ bytes of indexer keys
(8\,MB). The main attention does $k\cdot128\cdot(576+512)\approx2.9\times10^8$ multiply-adds against
$2048\cdot576\cdot2\approx2.4$\,MB of latent cache. Dense attention at the same length would need
$1.8\times10^{10}$ multiply-adds and 151\,MB of reads. The indexer is still $O(T)$ per token, but with a
constant small enough that the sparse layer is dominated by the $O(k)$ term up to very long contexts. (All
dimensions here are illustrative, chosen to be close to published MLA shapes.)
\end{napkinmath}


\begin{systemslens}[Sparsity has to align with tiles]
Kernels skip work at tile granularity, so a pattern is fast only if whole $B_r\times B_c$ tiles can be
skipped (\cref{fig:attention:patterns}). Sliding windows and block-sparse masks satisfy this; arbitrary
per-query top-$k$ selections do not, and need gather-based kernels whose memory access is less regular.
For decoding, the relevant question is bytes read per step: a window bounds them by $w$, and top-$k$
selection bounds the attention read at $k$ entries but still pays the indexer's full scan, at lower
precision and width.
\end{systemslens}

% =====================================================================
\section{Linear attention and its three forms}\label{sec:attnalt:linear}
\covers{CS336 L4}

Write causal attention for one head as $\vect{y}_i=\sum_{j\le i} a_{ij}\vect{v}_j$ with
$a_{ij}=\kappa(\vect{q}_i,\vect{k}_j)/\sum_{j'\le i}\kappa(\vect{q}_i,\vect{k}_{j'})$. Softmax attention
uses $\kappa(\vect{q},\vect{k})=\exp(\vect{q}^\top\vect{k}/\sqrt d)$. Linear attention (Katharopoulos et
al., 2020) uses a kernel that factorises through a feature map, $\kappa(\vect{q},\vect{k})=
\phi(\vect{q})^\top\phi(\vect{k})$; with $\phi(x)=\mathrm{elu}(x)+1$ the features stay positive. Then
\begin{equation}
\vect{y}_i^\top \;=\; \frac{\phi(\vect{q}_i)^\top\sum_{j\le i}\phi(\vect{k}_j)\vect{v}_j^\top}
{\phi(\vect{q}_i)^\top\sum_{j\le i}\phi(\vect{k}_j)}
\;=\;\frac{\phi(\vect{q}_i)^\top\mat{S}_i}{\phi(\vect{q}_i)^\top\vect{z}_i},
\qquad \mat{S}_i=\sum_{j\le i}\phi(\vect{k}_j)\vect{v}_j^\top .
\label{eq:attnalt:linear}
\end{equation}
The trick is associativity: $(\mat{Q}\mat{K}^\top)\mat{V}=\mat{Q}(\mat{K}^\top\mat{V})$ when nothing
nonlinear sits between the products. Below we write $\vect{q},\vect{k}$ for the feature-mapped vectors and
drop the normaliser, which is just the same computation with $\vect{v}$ replaced by 1 (many modern
variants drop it in favour of a normalisation layer on the output).

\begin{figure}[t]
\centering
\begin{tikzpicture}[
  m/.style={draw, fill=#1, minimum width=0.3cm, minimum height=0.3cm, inner sep=0pt},
  lbl/.style={font=\scriptsize}]
% left: (QK^T) V
\node[font=\small\bfseries] at (2.0,2.2) {$(\mat{Q}\mat{K}^\top)\,\mat{V}$: $O(T^2d)$};
\node[draw, fill=fcompute!25, minimum width=0.5cm, minimum height=1.6cm] (Q) at (0,0.6) {};
\node[lbl] at (0,-0.4) {$T\times d$};
\node[draw, fill=fcompute!25, minimum width=1.6cm, minimum height=0.5cm] (K) at (1.35,1.15) {};
\node[lbl] at (1.35,1.6) {$d\times T$};
\node[draw, fill=warnred!25, minimum width=1.6cm, minimum height=1.6cm] (A) at (1.35,0.05) {$T\times T$};
\node[draw, fill=fdata!25, minimum width=0.5cm, minimum height=1.6cm] (V) at (2.85,0.05) {};
\node[lbl] at (2.85,-0.95) {$T\times d$};
\node[lbl, text=warnred] at (1.35,-1.05) {grows as $T^2$};
% right: Q (K^T V)
\begin{scope}[xshift=6.6cm]
\node[font=\small\bfseries] at (1.6,2.2) {$\mat{Q}\,(\mat{K}^\top\mat{V})$: $O(Td^2)$};
\node[draw, fill=fcompute!25, minimum width=1.6cm, minimum height=0.5cm] (K2) at (0.9,1.2) {};
\node[lbl] at (0.9,1.65) {$d\times T$};
\node[draw, fill=fdata!25, minimum width=0.5cm, minimum height=1.6cm] (V2) at (2.15,1.2) {};
\node[lbl] at (2.85,1.2) {$T\times d$};
\node[draw, fill=keygreen!30, minimum width=0.5cm, minimum height=0.5cm] (S) at (0.9,-0.1) {$\mat{S}$};
\node[lbl, text=keygreen] at (0.9,-0.6) {$d\times d$ state};
\node[draw, fill=fcompute!25, minimum width=0.5cm, minimum height=1.6cm] (Q2) at (-0.4,-0.1) {};
\node[lbl] at (-0.4,-1.1) {$T\times d$};
\draw[-{Stealth}, thick, keygreen] (1.35,0.95) -- (1.0,0.15);
\end{scope}
\end{tikzpicture}
\caption{Without the softmax, the two products can be associated either way. Contracting over time
first produces a $d\times d$ state that does not grow with $T$; causality then turns this into a running
sum over time, \cref{eq:attnalt:recurrent}.}
\label{fig:attnalt:assoc}
\end{figure}

\paragraph{Form 1: parallel (quadratic).} Compute $\mat{Y}=(\mat{Q}\mat{K}^\top\odot\mat{M})\mat{V}$ with
the causal mask $\mat{M}$. This is ordinary masked attention without the softmax, so it uses matmuls and
tensor cores and costs $O(T^2d)$.

\paragraph{Form 2: recurrent.} Causality makes the state a running sum:
\begin{equation}
\mat{S}_t=\mat{S}_{t-1}+\vect{k}_t\vect{v}_t^\top,\qquad \vect{y}_t^\top=\vect{q}_t^\top\mat{S}_t .
\label{eq:attnalt:recurrent}
\end{equation}
This is an RNN whose hidden state is a $d_k\times d_v$ matrix. Each step costs $O(d_kd_v)$ regardless of
$t$, which makes decoding constant-time and constant-memory. Training this way is a sequential loop of
rank-1 updates, which is hopeless on a GPU.

\paragraph{Form 3: chunkwise.} Split the sequence into chunks of $C$ tokens. Within a chunk use the parallel
form; across chunks carry the state. For chunk $c$ with rows $\mat{Q}_c,\mat{K}_c,\mat{V}_c$ and incoming
state $\mat{S}$,
\begin{equation}
\mat{Y}_c=\underbrace{(\mat{Q}_c\mat{K}_c^\top\odot\mat{M})\mat{V}_c}_{\text{intra-chunk}}
+\underbrace{\mat{Q}_c\mat{S}}_{\text{inter-chunk}},
\qquad
\mat{S}\leftarrow\mat{S}+\mat{K}_c^\top\mat{V}_c .
\label{eq:attnalt:chunk}
\end{equation}
All three terms are matmuls of size $C\times d$ or $d\times d$, so tensor cores do the work, and the
sequential dependency is only across $T/C$ chunks. With $C=1$ this is the recurrent form and with $C=T$
the parallel form; in between, the cost is $O(TCd+Td^2)$. \Cref{fig:attnalt:chunk} sketches the dataflow,
and \cref{lst:attnalt:chunked} is the reference implementation, including the decay factors introduced
in the next section.

\begin{figure}[t]
\centering
\resizebox{0.95\linewidth}{!}{%
\begin{tikzpicture}[
  ch/.style={draw=sysblue, fill=sysbluebg, minimum width=2.2cm, minimum height=0.9cm, font=\scriptsize, align=center},
  st/.style={draw=keygreen, fill=keygreen!20, circle, minimum size=0.8cm, font=\scriptsize},
  arr/.style={-{Stealth[length=1.8mm]}, thick}]
\foreach \c/\x in {1/0,2/3.2,3/6.4,4/9.6} {
  \node[ch] (c\c) at (\x,0) {chunk \c\\$(\mat{Q}_c\mat{K}_c^\top\odot\mat{M})\mat{V}_c$};
  \node[st] (s\c) at (\x+1.6,-1.6) {$\mat{S}_{\c}$};
  \node[font=\scriptsize] (y\c) at (\x,1.2) {$\mat{Y}_\c$};
  \draw[arr] (c\c) -- (y\c);
}
\node[st] (s0) at (-1.6,-1.6) {$\mat{S}_0$};
\foreach \a/\b in {0/1,1/2,2/3,3/4} {
  \draw[arr, keygreen] (s\a) -- node[below, font=\tiny, text=black] {$+\mat{K}_\b^\top\mat{V}_\b$} (s\b);
  \draw[arr, keygreen, dashed] (s\a) to[out=60,in=-120] node[left, font=\tiny, text=black, pos=0.55] {$\mat{Q}_\b\mat{S}_\a$} (c\b.south);
}
\node[font=\scriptsize, align=left, text=black!70] at (13.6,0) {parallel within\\a chunk (tensor\\cores)};
\node[font=\scriptsize, align=left, text=black!70] at (13.6,-1.6) {sequential over\\$T/C$ chunks};
\end{tikzpicture}}
\caption{Chunkwise-parallel linear attention. Each chunk computes its intra-chunk term with a small
masked matmul, reads the incoming state for the inter-chunk term, and passes an updated state on.
Kernels such as those in the flash-linear-attention library keep $\mat{S}$ in on-chip memory between chunks.}
\label{fig:attnalt:chunk}
\end{figure}

\begin{lstlisting}[caption={Chunkwise linear attention with optional decay (\texttt{llmsys/linear\_attn.py}).},label={lst:attnalt:chunked}]
def linear_attention_chunked(q, k, v, gamma=None, normalize=False, chunk=16):
    B, H, T, dk = q.shape
    va = _augment(v, normalize)                   # append ones -> normaliser for free
    log_g = _log_gamma(gamma, B, H, T, q.device, q.dtype)
    S = q.new_zeros(B, H, dk, va.shape[-1])
    outs = []
    for s in range(0, T, chunk):
        e = min(s + chunk, T)
        qc, kc, vc, lg = q[..., s:e, :], k[..., s:e, :], va[..., s:e, :], log_g[..., s:e]
        cum = torch.cumsum(lg, dim=-1)            # log decay from chunk start to i
        intra = (torch.einsum("bhik,bhjk->bhij", qc, kc) * decay_mask(lg)) @ vc
        inter = torch.einsum("bhik,bhkv->bhiv", qc * cum.exp()[..., None], S)
        outs.append(intra + inter)
        total = cum[..., -1]                      # decay across the whole chunk
        w = (total[..., None] - cum).exp()        # decay from token j to chunk end
        S = total.exp()[..., None, None] * S \
            + torch.einsum("bhjk,bhjv->bhkv", kc * w[..., None], vc)
    return _maybe_normalize(torch.cat(outs, dim=2), normalize)
\end{lstlisting}

\begin{napkinmath}[When is linear attention cheaper?]
Per token and per layer, causal softmax attention does about $2Td$ multiply-adds for scores and outputs
($d=H\dhead$). Chunkwise linear attention does about $Cd$ for the intra-chunk part and $2d\dhead$ for
reading and updating the per-head states. With $\dhead=128$ and $C=64$ the linear form is cheaper once
$T\gtrsim (C+2\dhead)/2\approx160$. During decoding the comparison is in bytes: softmax reads $2T\dhead$
values per KV head, while linear attention reads a $\dhead\times\dhead$ state per head, equivalent to the KV
cache of only $\dhead/2=64$ tokens per head (or $\dhead$ tokens when the heads are not grouped).
\end{napkinmath}

\paragraph{Choosing $\phi$ and the normaliser.} The feature map decides how well $\phi(\vect{q})^\top\phi(\vect{k})$
imitates $\exp(\vect{q}^\top\vect{k})$. Performer (Choromanski et al., 2021) uses positive random features
whose expectation equals the softmax kernel exactly, at the cost of variance and an extra projection.
Simpler deterministic maps include $\mathrm{elu}+1$, ReLU, and low-order Taylor expansions of $\exp$
(as in the Based architecture), which approximate softmax well near zero and expand $d_k$ to $O(d_k^2)$
features. The division by $\phi(\vect{q})^\top\vect{z}_i$ in \cref{eq:attnalt:linear} is numerically awkward:
it can approach zero and it couples every output to a running sum. Most recent designs (RetNet, GLA,
Mamba-2, DeltaNet) drop it entirely and put a normalisation layer (RMSNorm or GroupNorm per head) on the
output instead, which is why our layer makes \code{normalize} optional.

\paragraph{Why plain linear attention underperforms.} Two weaknesses kept linear attention a curiosity
for years. First, it never forgets: $\mat{S}$ is a sum of everything ever seen, so old content blurs new
retrievals. Second, a $d_k\times d_v$ matrix can store only about $d_k$ key--value associations without
interference, while softmax attention can retrieve any of $T$ tokens sharply. The next two sections fix the
first problem with decay and gating and attack the second with the delta rule. The capacity limit
remains, which is the reason hybrids keep some softmax layers.

% =====================================================================
\section{Decay and gating: RetNet, GLA, and Mamba-2}\label{sec:attnalt:gating}
\covers{CS336 L4}

Add a forgetting factor to \cref{eq:attnalt:recurrent}:
\begin{equation}
\mat{S}_t=\gamma_t\,\mat{S}_{t-1}+\vect{k}_t\vect{v}_t^\top,\qquad\vect{y}_t^\top=\vect{q}_t^\top\mat{S}_t,
\qquad 0<\gamma_t\le1 .
\label{eq:attnalt:gated}
\end{equation}
Unrolling gives $\vect{y}_i^\top=\sum_{j\le i}\big(\prod_{l=j+1}^{i}\gamma_l\big)(\vect{q}_i^\top\vect{k}_j)\vect{v}_j^\top$:
the parallel form is masked attention with a \emph{decay mask} $D_{ij}=\prod_{l=j+1}^i\gamma_l$ in place of the
0/1 causal mask. The chunked form needs only cumulative products of $\gamma$ within each chunk, computed
stably as cumulative sums of $\log\gamma$ (\cref{lst:attnalt:chunked}). The variants differ in where
$\gamma$ comes from:
\begin{itemize}
\item \textbf{RetNet} (Sun et al., 2023): a fixed per-head constant, $\gamma_h=1-2^{-5-h}$, so different
  heads have different memory horizons.
\item \textbf{Gated linear attention} (GLA; Yang et al., 2024) and \textbf{Mamba-2} (Dao and Gu, 2024):
  $\gamma_t$ is computed from the input, so the model decides when to forget, for example at a document
  boundary. Mamba-2 uses a scalar per head and per step. Its ``state space duality'' (SSD) shows that the
  selective SSM of \cref{sec:attnalt:ssm} with scalar decay is exactly this structured masked attention,
  and derives the chunked algorithm from that equivalence.
\end{itemize}

\begin{lstlisting}[caption={A gated linear-attention layer with a constant-memory decoding step.},label={lst:attnalt:layer}]
class LinearAttention(nn.Module):
    def __init__(self, d_model, n_heads, mode="gated", chunk=16):
        super().__init__()
        self.H, self.dh, self.chunk = n_heads, d_model // n_heads, chunk
        self.qkv = nn.Linear(d_model, 3 * d_model, bias=False)
        self.out = nn.Linear(d_model, d_model, bias=False)
        self.gate = nn.Linear(d_model, n_heads)          # one forget gate per head/step
        nn.init.constant_(self.gate.bias, 3.0)           # sigmoid(3) = 0.95: slow decay

    def forward(self, x):                                # training: chunkwise parallel
        q, k, v, g = self._proj(x)                       # g = sigmoid(gate(x)), (B,H,T)
        y = linear_attention_chunked(q, k, v, g, normalize=True, chunk=self.chunk)
        return self.out(merge_heads(y))

    @torch.no_grad()
    def step(self, x_t, state=None):                     # decoding: O(1) per token
        q, k, v, g = self._proj(x_t)
        y, state = linear_attention_recurrent(q, k, v, g, True, state)
        return self.out(merge_heads(y)), state
\end{lstlisting}

The test suite of Lab~8 checks that the parallel, recurrent and chunked forms agree to $10^{-10}$ in
float64 for every chunk size and every kind of decay, and that \code{step} reproduces \code{forward}
token by token. This kind of equivalence test is how you should develop any new sequence mixer:
write the slow, obviously correct form first, then the fast one, and compare them.

\begin{systemslens}[Choosing the chunk size]
The chunk size $C$ trades the $O(TCd)$ intra-chunk work, which is matmul-friendly, against $T/C$ sequential
state updates, each reading and writing a $d_k\times d_v$ state per head. On GPUs $C=64$ or 128 matches
tensor-core tiles, and production kernels keep the state in shared memory or registers across chunks so
that only $\mat{Q},\mat{K},\mat{V},\mat{Y}$ touch HBM. The decay mask within a chunk involves
$\exp$ of differences of cumulative log-gates, so kernels compute it in fp32. Products of many $\gamma$s
underflow quickly, and a careless implementation divides by them, which produces NaNs. The reference
code only ever exponentiates \emph{non-positive} differences.
\end{systemslens}

% =====================================================================
\section{The delta rule and gated DeltaNet}\label{sec:attnalt:delta}
\covers{CS336 L4}

Viewed as memory, $\mat{S}$ is an associative store: writing $\vect{k}\vect{v}^\top$ files value
$\vect{v}$ under key $\vect{k}$, and reading with $\vect{q}=\vect{k}$ returns $\mat{S}^\top\vect{k}$. With
additive writes, re-filing a key \emph{adds} the new value to the old one. The delta rule (Widrow--Hoff,
adopted for fast-weight transformers by Schlag et al., 2021) instead writes only the \emph{error}:
\begin{equation}
\mat{S}_t=\mat{S}_{t-1}+\beta_t\,\vect{k}_t\big(\vect{v}_t-\mat{S}_{t-1}^\top\vect{k}_t\big)^\top
=\big(\mat{I}-\beta_t\vect{k}_t\vect{k}_t^\top\big)\mat{S}_{t-1}+\beta_t\vect{k}_t\vect{v}_t^\top .
\label{eq:attnalt:delta}
\end{equation}
With a unit-norm key and $\beta_t=1$, the update replaces whatever was stored under $\vect{k}_t$ by
$\vect{v}_t$ exactly: it is one step of gradient descent with learning rate $\beta_t$ on the regression loss
$\frac12\|\mat{S}^\top\vect{k}_t-\vect{v}_t\|^2$. This view connects linear recurrent models to
\emph{test-time training}, where a model's state is a set of weights updated by an inner optimiser as it
reads. Gated DeltaNet (Yang et al., 2025) multiplies the whole update by a forget gate $\gamma_t$, combining
the two mechanisms:
\[
\mat{S}_t=\gamma_t\big(\mat{I}-\beta_t\vect{k}_t\vect{k}_t^\top\big)\mat{S}_{t-1}+\gamma_t\beta_t\vect{k}_t\vect{v}_t^\top .
\]
Qwen's recent ``Next'' models are reported to interleave gated-DeltaNet layers with full attention at
roughly 3:1.

The transition matrix $\mat{I}-\beta\vect{k}\vect{k}^\top$ is no longer a scalar, so the simple decay-mask
trick does not apply. The chunked algorithm instead represents a product of such Householder-like matrices
compactly with the WY representation from numerical linear algebra, which again reduces the work to
matmuls (Yang et al., 2024). The lab implements the recurrent step, and the test demonstrates the
overwrite property directly.

\begin{lstlisting}[caption={One delta-rule step (Lab 8, Task 4).},label={lst:attnalt:delta}]
def delta_rule_step(S, q_t, k_t, v_t, beta_t, g_t=None):
    pred = torch.einsum("bhk,bhkv->bhv", k_t, S)          # what memory returns for k_t
    S = S + beta_t[..., None, None] * k_t[..., :, None] * (v_t - pred)[..., None, :]
    if g_t is not None:
        S = g_t[..., None, None] * S                      # gated DeltaNet: forget
    return torch.einsum("bhk,bhkv->bhv", q_t, S), S
\end{lstlisting}

% =====================================================================
\section{State-space models and parallel scans}\label{sec:attnalt:ssm}
\covers{CS336 L4}

State-space models arrive at the same place from signal processing. A continuous linear system maps an
input signal $x(t)$ to an output $y(t)$ through a hidden state $h(t)\in\R^N$:
\[
h'(t)=\mat{A}h(t)+\mat{B}x(t),\qquad y(t)=\mat{C}h(t)+Dx(t).
\]
Discretising with step size $\Delta$ (zero-order hold for $\mat{A}$, a first-order approximation for
$\mat{B}$) gives the recurrence
\begin{equation}
h_t=\bar{\mat{A}}_t\,h_{t-1}+\bar{\mat{B}}_t\,x_t,\qquad \bar{\mat{A}}_t=\exp(\Delta_t\mat{A}),\quad
\bar{\mat{B}}_t=\Delta_t\mat{B}_t,\qquad y_t=\mat{C}_th_t+Dx_t .
\label{eq:attnalt:ssm}
\end{equation}

\paragraph{S4 to Mamba.} In S4 (Gu et al., 2022) the parameters are time-invariant, so the whole map is a
long convolution that FFTs compute in parallel, with careful structure on $\mat{A}$ for long memory. Such
LTI models cannot \emph{select}: they treat every token identically, so they cannot ignore filler or
latch onto a key. Mamba (Gu and Dao, 2023) makes $\Delta_t$, $\mat{B}_t$ and $\mat{C}_t$ functions of the
input. A large $\Delta_t$ resets the state toward the current input and a small one preserves the past, so
the model can choose what to remember, much like the gates above. With diagonal $\mat{A}$, each of the $E$
channels runs $N$ independent scalar recurrences.

\paragraph{The scan.} Selectivity destroys the convolution view, but \cref{eq:attnalt:ssm} is still a
\emph{first-order linear recurrence} $h_t=a_th_{t-1}+b_t$, and such recurrences compose associatively.
Represent step $t$ by the pair $(a_t,b_t)$ meaning ``$h\mapsto a_th+b_t$''. Applying step $s$ then step $t$
gives $h\mapsto a_t(a_sh+b_s)+b_t$, so
\begin{equation}
(a_s,b_s)\bullet(a_t,b_t)=(a_sa_t,\;a_tb_s+b_t),
\label{eq:attnalt:scanop}
\end{equation}
and $\bullet$ is associative. Any associative operator admits a \emph{parallel prefix scan}: all $T$ prefix
results in $O(\log T)$ parallel rounds (Hillis--Steele, $O(T\log T)$ work) or with $O(T)$ work in two sweeps
(Blelloch). \Cref{fig:attnalt:scan} shows the Hillis--Steele pattern.

\begin{figure}[t]
\centering
\begin{tikzpicture}[x=1.25cm, y=0.95cm,
  n/.style={circle, draw, minimum size=0.55cm, inner sep=0pt, font=\tiny},
  e/.style={-{Stealth[length=1.4mm]}, sysblue, thick}]
\foreach \t in {0,...,7} {
  \node[n, fill=fdata!20] (r0\t) at (\t,0) {$\t$};
  \node[n, fill=white] (r1\t) at (\t,-1) {};
  \node[n, fill=white] (r2\t) at (\t,-2) {};
  \node[n, fill=keygreen!25] (r3\t) at (\t,-3) {$0{:}\t$};
}
\foreach \t in {0,...,7} { \draw[black!30] (r0\t)--(r1\t); \draw[black!30] (r1\t)--(r2\t); \draw[black!30] (r2\t)--(r3\t); }
\foreach \t in {1,...,7} { \pgfmathtruncatemacro{\s}{\t-1} \draw[e] (r0\s) -- (r1\t); }
\foreach \t in {2,...,7} { \pgfmathtruncatemacro{\s}{\t-2} \draw[e] (r1\s) -- (r2\t); }
\foreach \t in {4,...,7} { \pgfmathtruncatemacro{\s}{\t-4} \draw[e] (r2\s) -- (r3\t); }
\node[font=\scriptsize, anchor=west] at (7.5,-0.5) {offset 1};
\node[font=\scriptsize, anchor=west] at (7.5,-1.5) {offset 2};
\node[font=\scriptsize, anchor=west] at (7.5,-2.5) {offset 4};
\end{tikzpicture}
\caption{Hillis--Steele inclusive scan over 8 steps. In round $r$, element $t$ combines with element
$t-2^r$ using \cref{eq:attnalt:scanop}; after $\lceil\log_2 T\rceil$ rounds element $t$ holds the
composition of steps $0..t$, i.e.\ $h_t$.}
\label{fig:attnalt:scan}
\end{figure}

\begin{lstlisting}[caption={Parallel scan for $h_t=a_th_{t-1}+b_t$ (\texttt{llmsys/ssm.py}).},label={lst:attnalt:scan}]
def scan_parallel(a, b):
    a, b = a.clone(), b.clone()
    T, step = a.shape[1], 1
    while step < T:                               # ceil(log2 T) rounds
        a_prev = torch.ones_like(a);  a_prev[:, step:] = a[:, :-step]
        b_prev = torch.zeros_like(b); b_prev[:, step:] = b[:, :-step]
        b = a * b_prev + b                        # (a_prev, b_prev) . (a, b)
        a = a * a_prev
        step *= 2
    return b                                      # b[:, t] == h_t
\end{lstlisting}

\paragraph{The Mamba block.} \code{MambaBlock} in \code{llmsys/ssm.py} follows the published design: an input
projection to two streams, a short causal depthwise convolution (width 4) on one, SiLU, the selective SSM,
multiplicative gating by SiLU of the other stream, and an output projection. Its \code{step} method
carries two pieces of state, the last three conv inputs and $h\in\R^{E\times N}$, and reproduces the
parallel forward pass token by token.


\paragraph{Discretisation in detail.} For a single diagonal entry $a<0$ of $\mat{A}$ and an input held
constant over an interval of length $\Delta$ (zero-order hold), the exact solution of $h'=ah+bx$ is
\[
h(t+\Delta)=e^{\Delta a}h(t)+\frac{e^{\Delta a}-1}{a}\,b\,x(t).
\]
The first factor gives $\bar{a}=\exp(\Delta a)\in(0,1)$, a decay just like $\gamma$ in
\cref{eq:attnalt:gated}. For small $\Delta a$, $(e^{\Delta a}-1)/a\approx\Delta$, which is the simpler
$\bar{b}=\Delta b$ that Mamba's implementation uses. Because $a<0$ and $\Delta>0$ (a softplus output), the
decay always lies in $(0,1)$ and the recurrence is stable by construction, which is why $\mat{A}$ is
parameterised as $-\exp(\mat{A}_{\log})$. When $\mat{A},\mat{B},\mat{C},\Delta$ do not depend on $t$ (S4),
unrolling gives $y=\mat{K}*x$ with the convolution kernel
$\mat{K}=(\mat{C}\bar{\mat{B}},\,\mat{C}\bar{\mat{A}}\bar{\mat{B}},\,\mat{C}\bar{\mat{A}}^2\bar{\mat{B}},\ldots)$,
which FFTs compute in $O(T\log T)$. Selectivity makes $\bar{\mat{A}}_t$ time-varying and removes this option.

\begin{figure}[t]
\centering
\begin{tikzpicture}[
  b/.style={draw, rounded corners=2pt, minimum width=1.7cm, minimum height=0.55cm, font=\scriptsize, align=center},
  arr/.style={-{Stealth[length=1.6mm]}, thick}]
\node[b, fill=fcompute!20] (in) at (0,0) {in\_proj\\$d\to 2E$};
\node[b, fill=fdata!20] (conv) at (-1.5,1.2) {causal conv1d\\(width 4)};
\node[b, fill=fdata!20] (act) at (-1.5,2.2) {SiLU};
\node[b, fill=keygreen!25] (ssm) at (-1.5,3.3) {selective SSM\\$\Delta_t,\mat{B}_t,\mat{C}_t$ from input};
\node[b, fill=fdata!20] (gate) at (1.5,2.2) {SiLU};
\node[circle, draw, inner sep=1pt, font=\scriptsize] (mul) at (0,4.4) {$\times$};
\node[b, fill=fcompute!20] (out) at (0,5.3) {out\_proj\\$E\to d$};
\draw[arr] (0,-0.8) -- (in);
\draw[arr] (in) -| node[pos=0.75, left, font=\tiny] {$x$} (conv);
\draw[arr] (in) -| node[pos=0.75, right, font=\tiny] {$z$} (gate);
\draw[arr] (conv) -- (act); \draw[arr] (act) -- (ssm);
\draw[arr] (ssm) |- (mul); \draw[arr] (gate) |- (mul);
\draw[arr] (mul) -- (out); \draw[arr] (out) -- (0,6.0);
\node[font=\scriptsize, align=left, text=black!70, anchor=west] at (2.6,1.2) {decode state:\\last 3 conv inputs};
\node[font=\scriptsize, align=left, text=black!70, anchor=west] at (2.6,3.3) {decode state:\\$h\in\R^{E\times N}$};
\end{tikzpicture}
\caption{The Mamba block as implemented in \code{llmsys/ssm.py}. Two pieces of state are carried during
decoding, both of fixed size. The short convolution gives each token a little local context before the
recurrence; the gate plays the role of a GLU.}
\label{fig:attnalt:mamba}
\end{figure}

\begin{lstlisting}[caption={The selective scan and the Mamba block's decode step (\texttt{llmsys/ssm.py}, abridged).},label={lst:attnalt:mamba}]
def selective_scan(x, dt, A, Bm, Cm, D, parallel=True, h0=None):
    dA = torch.exp(dt[..., None] * A)                        # (B,T,E,N) decay
    dBx = dt[..., None] * Bm[:, :, None, :] * x[..., None]   # (B,T,E,N) input
    h = scan_parallel(dA, dBx) if parallel and h0 is None \
        else scan_sequential(dA, dBx, h0)
    y = torch.einsum("bten,btn->bte", h, Cm) + D * x
    return y, h[:, -1]

@torch.no_grad()
def step(self, u_t, state):                                  # u_t: (B, 1, d_model)
    conv_s, h = state                                        # (B,E,3), (B,E,N)
    x, z = self.in_proj(u_t).chunk(2, dim=-1)
    window = torch.cat([conv_s, x.transpose(1, 2)], dim=-1)  # last 4 inputs
    xc = F.silu((window * self.conv.weight.squeeze(1)).sum(-1) + self.conv.bias)[:, None]
    dt, Bm, Cm = self._ssm_params(xc)
    y, h = selective_scan(xc, dt, -torch.exp(self.A_log), Bm, Cm, self.D,
                          parallel=False, h0=h)
    return self.out_proj(y * F.silu(z)), (window[..., 1:], h)
\end{lstlisting}

\begin{systemslens}[Hardware-aware scans]
The state $h$ has $E\cdot N$ entries per token position; materialising it for every $t$ in HBM, as our
reference does, costs $T\cdot E\cdot N$ values, far more than the activations themselves. Mamba's kernel
avoids this: it loads $\Delta,\mat{A},\mat{B},\mat{C},x$ for a block into SRAM, runs the discretisation and
the scan there, writes only $y$, and \emph{recomputes} the states in the backward pass. This is the same
IO-aware philosophy as FlashAttention (\cref{sec:attention:flash}). Mamba-2's SSD reformulation goes
further by turning most of the work into chunked matmuls (\cref{eq:attnalt:chunk} with decay), because a
scan of elementwise operations cannot use tensor cores and a matmul can.
\end{systemslens}

% =====================================================================
\section{Hybrid architectures}\label{sec:attnalt:hybrid}
\covers{CS336 L4}

Fixed-size states are cheap but lossy; softmax attention is expensive but can retrieve any token exactly.
Pure recurrent models struggle on tasks that need verbatim recall from far back, such as copying,
associative recall and needle-in-a-haystack retrieval. Hybrids keep a few attention layers for retrieval and
use cheap mixers everywhere else. Reported designs include Jamba (one attention layer per seven Mamba layers,
with MoE MLPs), MiniMax-01 (seven ``lightning'' linear-attention layers per softmax layer), NVIDIA's
Nemotron-H family (mostly Mamba-2 with a few attention layers), and Qwen's ``Next'' models (gated-DeltaNet
and attention at about 3:1). Controlled published ablations are still scarce, but a consistent finding is
that a small fraction of attention layers recovers most of the gap to a full-attention model.

\begin{figure}[t]
\centering
\begin{tikzpicture}
\begin{axis}[width=0.78\linewidth, height=5.6cm, xmode=log, ymode=log, log basis x=2,
  xlabel={context length $T$ (tokens)}, ylabel={per-sequence cache (bytes)},
  xmin=1024, xmax=1048576, grid=both, grid style={black!8},
  legend pos=north west, legend style={font=\scriptsize}, tick label style={font=\scriptsize}]
% 36 layers, 8 KV heads, 32 query heads, dh=128, bf16
\addplot[warnred, thick, domain=1024:1048576, samples=20] {2*36*8*128*x*2};
\addlegendentry{full attention (GQA 8)}
\addplot[napkin, thick, dashed, domain=1024:1048576, samples=40] {2*36*8*128*min(x,4096)*2};
\addlegendentry{sliding window $w=4096$}
\addplot[sysblue, thick, domain=1024:1048576, samples=20] {2*9*8*128*x*2 + 27*32*128*128*2};
\addlegendentry{hybrid 3 linear : 1 full}
\addplot[keygreen, thick, domain=1024:1048576, samples=20] {36*32*128*128*2};
\addlegendentry{all linear / SSM state}
\end{axis}
\end{tikzpicture}
\caption{Decoding memory for one sequence of a Qwen3-8B-shaped stack (36 layers, 32 query heads, 8 KV
heads, $\dhead=128$, bf16) as the sequence mixer changes. The hybrid's cache still grows linearly but four
times more slowly; the pure recurrent state is constant at 37.7\,MB, the KV cache of 256 tokens.}
\label{fig:attnalt:cache}
\end{figure}

\paragraph{Why some attention is needed: associative recall.} The cleanest diagnostic is
\emph{multi-query associative recall} (MQAR; Arora et al., 2023). The input lists key--value pairs, and later
queries repeat some of the keys; the model must output the matching values (\cref{fig:attnalt:mqar}).
Softmax attention solves it with one layer at any number of pairs, because it can look up the exact earlier
token. A fixed-size state must \emph{compress} all pairs into $d_k\times d_v$ numbers, so accuracy falls
once the number of pairs exceeds the state's capacity. Accuracy on MQAR across state sizes tracked
language-modelling quality gaps between attention and recurrent models remarkably well, and it is the
quickest test to run on a new mixer.

\begin{figure}[t]
\centering
\begin{tikzpicture}[
  tok/.style={draw, minimum width=0.55cm, minimum height=0.5cm, font=\scriptsize\ttfamily, inner sep=1pt},
  kk/.style={tok, fill=sysblue!20}, vv/.style={tok, fill=keygreen!25},
  qq/.style={tok, fill=napkin!25}, aa/.style={tok, fill=keygreen!45, draw=keygreen, thick}]
\foreach \x/\k/\v in {0/A/4,1.2/C/9,2.4/F/1,3.6/B/7} {
  \node[kk] at (\x,0) {\k}; \node[vv] at (\x+0.55,0) {\v};
}
\node[font=\scriptsize] at (5.0,0) {$\cdots$};
\foreach \x/\q/\a in {5.8/F/1,7.0/A/4} {
  \node[qq] at (\x,0) {\q}; \node[aa] at (\x+0.55,0) {\a};
}
\draw[-{Stealth}, keygreen, thick] (6.35,0.3) to[out=140,in=40] (2.95,0.3);
\draw[-{Stealth}, keygreen, thick] (7.55,0.3) to[out=150,in=30] (0.55,0.3);
\node[font=\scriptsize, align=left, anchor=west] at (8.1,0) {predict the value\\filed under the key};
\node[font=\scriptsize, text=black!70] at (1.8,-0.6) {key--value pairs};
\node[font=\scriptsize, text=black!70] at (6.45,-0.6) {queries};
\end{tikzpicture}
\caption{Multi-query associative recall. Attention can point back at the exact earlier token (arrows); a
recurrent model must have stored the pair in its fixed-size state.}
\label{fig:attnalt:mqar}
\end{figure}

\begin{systemslens}[Training recurrent layers on long sequences across GPUs]
For softmax attention, context parallelism shards the sequence and circulates KV blocks around a ring
(\cref{ch:distributed}). Recurrent layers have a simpler but more \emph{sequential} structure: device $r$ needs
the state at the end of device $r-1$'s shard. One approach computes each shard's local outputs and its
local state contribution in parallel and then fixes up the cross-shard term once the incoming states are
known. Because the state transition is linear, the incoming states can themselves be obtained with a scan
over devices: each device sends a single $d_k\times d_v$ matrix (plus a decay scalar) per head, which is far
less traffic than shipping KV blocks.
\end{systemslens}

\begin{systemslens}[Serving a hybrid]
Hybrids complicate the serving engine of \cref{ch:serving}. Attention layers want a paged KV cache, while
recurrent layers want one fixed-size state slot per sequence, which needs no paging but must be copied
for beam search or parallel sampling and cannot be partially reused. Prefix caching becomes harder because
a recurrent state summarises an \emph{entire} prefix: to reuse a shared prompt you must have checkpointed
the state exactly at the prefix boundary, not just kept per-token KV blocks. Speculative decoding needs the
state to roll back on rejection, so engines snapshot states per draft position. Chunked prefill works
naturally, since the chunkwise algorithm is chunked prefill. These are solvable problems, but they are
why support for new mixers in inference engines usually lags the model releases.
\end{systemslens}


% =====================================================================
\section{Setting up a hybrid model end to end}\label{sec:attnalt:setup}

New sequence mixers fail in practice far more often through plumbing than through modelling. The
workflow that keeps them correct is the same one used throughout this chapter:
\begin{enumerate}
\item \textbf{Write the obviously correct form first}: the recurrent loop, or the quadratic masked form,
  in float64, with no cleverness.
\item \textbf{Write the fast form} (chunked, scan-based, or a fused kernel) and compare it against the slow
  one on random inputs, over many shapes, chunk sizes and gate regimes, including gates near 0 and near 1.
\item \textbf{Write the decode step} and check that feeding a sequence one token at a time through
  \code{step} reproduces \code{forward} at every position. This catches off-by-one errors in conv windows,
  RoPE offsets and state updates.
\item \textbf{Only then integrate} the mixer into a model, a training loop and a serving path.
\end{enumerate}

\begin{lstlisting}[caption={The equivalence-test pattern used for every mixer in this chapter (from Lab 8).},label={lst:attnalt:test}]
@pytest.mark.parametrize("gamma", [None, per_head_decay, per_step_gates])
def test_three_forms_agree(gamma):
    q, k, v = random_qkv(dtype=torch.float64)            # float64: compare to 1e-10
    ref = linear_attention_quadratic(q, k, v, gamma)      # slow, obviously correct
    rec, _ = linear_attention_recurrent(q, k, v, gamma)   # decode path
    assert torch.allclose(ref, rec, atol=1e-10)
    for chunk in (1, 4, 7, 64):                           # include ragged last chunks
        assert torch.allclose(ref, linear_attention_chunked(q, k, v, gamma, chunk=chunk),
                              atol=1e-10)

def test_step_matches_forward():
    layer = LinearAttention(16, 2).eval()
    x = torch.randn(2, 20, 16)
    full, state, outs = layer(x), None, []
    for t in range(x.shape[1]):                           # one token at a time
        y, state = layer.step(x[:, t:t + 1], state)
        outs.append(y)
    assert torch.allclose(torch.cat(outs, 1), full, atol=1e-5)
\end{lstlisting}
Two details are easy to miss. Test in float64 so that a tolerance of $10^{-10}$ separates bugs from rounding,
and include chunk sizes that do not divide $T$, because ragged last chunks are where index arithmetic breaks.


\code{HybridLM} in \code{llmsys/linear\_attn.py} is a small but complete example of step 4. A pattern string
such as \code{"LLLA"} chooses the mixer for each layer: gated linear attention (L), Mamba (M) or softmax
attention with RoPE and a KV cache (A). Each layer is a pre-norm residual block with an MLP. The same
weights serve a parallel \code{forward} for training and a \code{decode\_step} that threads a
heterogeneous list of per-layer states through the network, as in \cref{lst:attnalt:hybrid}.

\begin{lstlisting}[caption={Decoding a hybrid model: one state object per layer, of a type that depends on the layer.},label={lst:attnalt:hybrid}]
@torch.no_grad()
def decode_step(self, ids_t, states):          # ids_t: (B, 1); states: list per layer
    x = self.emb(ids_t)
    new = []
    for p, mix, n1, n2, mlp, st in zip(self.pattern, self.mixers, self.n1,
                                       self.n2, self.mlps, states):
        h = n1(x)
        if p == "A":                            # KV cache: grows by one entry per step
            y, st = mix(h, st)
        elif p == "L":                          # (B, H, dk, dv+1) matrix state, fixed
            y, st = mix.step(h, st)
        else:                                   # Mamba: (conv window, SSM state), fixed
            st = st if st is not None else mix.init_state(ids_t.shape[0], ids_t.device)
            y, st = mix.step(h, st)
        new.append(st)
        x = x + y
        x = x + mlp(n2(x))
    return self.head(self.norm(x)), new

model = HybridLM(pattern="LMLA").eval()
states = [None] * len(model.pattern)
for t in range(ids.shape[1]):                   # equals model(ids) position by position
    logits_t, states = model.decode_step(ids[:, t:t + 1], states)
print(model.cache_numel(states))                # only the 'A' entry grows with t
\end{lstlisting}

\begin{figure}[t]
\centering
\begin{tikzpicture}[
  ly/.style={draw, rounded corners=1.5pt, minimum width=2.3cm, minimum height=0.55cm, font=\scriptsize},
  cache/.style={draw, minimum height=0.45cm, font=\tiny, anchor=west}]
\foreach \i/\t/\c in {0/L/keygreen,1/M/keygreen,2/L/keygreen,3/A/warnred} {
  \node[ly, fill=\c!18, draw=\c] (l\i) at (0,\i*0.85) {layer \i: \t};
}
\foreach \i in {0,2} { \node[cache, fill=keygreen!25, minimum width=0.9cm] at (1.6,\i*0.85) {$\mat{S}$: $H\!\times\! d_k\!\times\! d_v$}; }
\node[cache, fill=keygreen!25, minimum width=0.9cm] at (1.6,0.85) {(conv, $h$): $E\!\times\!(3+N)$};
\node[cache, fill=warnred!25, minimum width=4.6cm] at (1.6,2.55) {K, V: $2\times H\times t\times d_h$ (grows)};
\draw[-{Stealth}, thick, black!50] (6.5,2.55) -- (7.3,2.55) node[right, font=\scriptsize] {$+1$ per token};
\node[font=\scriptsize, align=center] at (0,-0.75) {pattern \texttt{"LMLA"}};
\end{tikzpicture}
\caption{The decode-time state of a hybrid stack. Only the attention layers hold memory that grows with the
generated length; the serving engine must manage paged blocks for them and fixed slots for the others.}
\label{fig:attnalt:hybridstate}
\end{figure}

To train it, use any loop from \cref{ch:pretrain} or \code{llmsys.ablate.train\_lm}'s structure: the model's
\code{forward(ids)} returns logits like every other model in the book. To profile it, time
\code{forward} at increasing $T$ for patterns \code{"AAAA"}, \code{"LLLA"} and \code{"LLLL"}: the
quadratic term appears only in proportion to the number of A layers. To serve it, the engine of
\cref{ch:serving} needs a second kind of cache slot, as discussed in the Systems Lens above. Lab~8's
stretch goals walk through each of these steps. \Cref{fig:attnalt:flops} shows what the profile should look like analytically.

\begin{figure}[t]
\centering
\begin{tikzpicture}
\begin{axis}[width=0.78\linewidth, height=5.4cm, xmode=log, ymode=log, log basis x=2,
  xlabel={position $T$ in the sequence}, ylabel={mixer MACs per token per layer},
  xmin=256, xmax=1048576, grid=both, grid style={black!8}, legend pos=north west,
  legend style={font=\scriptsize}, tick label style={font=\scriptsize}]
% d = 4096, dh = 128, C = 64 ; softmax ~ 2 T d ; linear ~ C d + 2 d dh
\addplot[warnred, thick, domain=256:1048576, samples=30] {2*x*4096}; \addlegendentry{softmax attention}
\addplot[sysblue, thick, domain=256:1048576, samples=30] {0.25*2*x*4096 + 0.75*(64*4096 + 2*4096*128)}; \addlegendentry{hybrid, 1 in 4 layers softmax}
\addplot[keygreen, thick, domain=256:1048576, samples=30] {64*4096 + 2*4096*128}; \addlegendentry{chunked linear ($C=64$)}
\addplot[black!50, densely dotted, domain=256:1048576, samples=2] {3*4096*12288}; \addlegendentry{SwiGLU MLP (for scale)}
\end{axis}
\end{tikzpicture}
\caption{Analytic per-token cost of the sequence mixer for a $d=4096$, $\dhead=128$ layer, averaged over a
causal prefix of length $T$, with the MLP's fixed cost for comparison. Softmax attention overtakes the MLP
near $T=18$k tokens; a 1:3 hybrid pushes that point four times further out; a pure linear mixer stays
constant and small.}
\label{fig:attnalt:flops}
\end{figure}

\begin{deepdive}[Test-time training and learned memories]
The delta rule is one step of SGD on a regression loss whose ``weights'' are the state. Generalising that
view gives \emph{test-time training} layers, in which the state is the parameters of a small model (linear or
an MLP) trained online, as the sequence is read, by an inner optimiser with its own learning rate,
momentum or weight decay. Recent proposals in this family (TTT layers; Titans, with a neural long-term
memory) differ in the inner model and update rule, and they inherit the systems questions of this chapter: how
to chunk the inner updates so that tensor cores do the work, and how large a state each sequence must carry
at decode time.
\end{deepdive}


% =====================================================================
\section{Multi-head latent attention, revisited}\label{sec:attnalt:mla}
\covers{CS336 L4}

\Cref{ch:modern} introduced MLA as ``caching a compressed latent''. Here are the details that make it work.
For each token, MLA projects the hidden state down to a small latent $\vect{c}_t=\mat{W}^{DKV}\vect{h}_t\in\R^{d_c}$
and caches only that. Keys and values for head $i$ are reconstructed by up-projections:
$\vect{k}_t^{(i)}=\mat{W}^{UK}_i\vect{c}_t$, $\vect{v}_t^{(i)}=\mat{W}^{UV}_i\vect{c}_t$.

\paragraph{Absorption.} At decode time nobody actually materialises the keys. The score
$\vect{q}^{(i)\top}\vect{k}_t^{(i)}=(\mat{W}^{UK\top}_i\vect{q}^{(i)})^\top\vect{c}_t$ can be computed
against the latent directly by folding $\mat{W}^{UK}_i$ into the query projection; similarly,
$\mat{W}^{UV}_i$ folds into the output projection. Attention then runs over $d_c$-dimensional cached
vectors shared by all heads, much like MQA with a wide head.

\paragraph{The RoPE problem.} Absorption requires the score to be bilinear in fixed matrices, but RoPE
inserts a position-dependent rotation $\mat{R}_{t}$ between $\vect{q}$ and $\vect{k}$, which cannot be
folded into a fixed $\mat{W}^{UK}$. DeepSeek-V2's fix is \emph{decoupled RoPE}: a small extra key component
$\vect{k}^{R}_t\in\R^{d_R}$, shared across heads and rotated normally, is cached alongside $\vect{c}_t$, and
the score becomes the sum of an absorbed no-position part and a small rotated part.

\begin{napkinmath}[MLA cache size]
DeepSeek-V2's published configuration uses 128 heads of key dimension 128 (plus 64 decoupled RoPE dimensions),
$d_c=512$ and $d_R=64$. Plain MHA would cache $2\times128\times128=32{,}768$ values per token per layer; MLA
caches $512+64=576$, which is 57 times less and comparable to GQA with only about two KV heads. Because the
absorbed decode attention reads one 576-wide vector per token for all 128 heads, its arithmetic intensity
per byte of cache is much higher than GQA's, which makes long-context decoding far less memory-bound.
\end{napkinmath}

\paragraph{Multi-token prediction.} DeepSeek-V3 also trains lightweight extra heads that predict one token
further ahead (MTP). It is a training signal, and at inference time the extra head can serve as a built-in
draft model for speculative decoding (\cref{ch:efficient}).

% =====================================================================
\section{Sparse experts, revisited}\label{sec:attnalt:moe}
\covers{CS336 L4}

\Cref{sec:modern:moe} implemented a top-$k$ MoE layer with an auxiliary balancing loss and expert
parallelism. Current MoE designs add several refinements.

\paragraph{Fine-grained and shared experts.} DeepSeekMoE splits each expert into $m$ smaller ones and routes to
$m$ times as many, which keeps FLOPs fixed while making the combinatorics of expert selection far richer.
It also adds a few \emph{shared} experts that every token uses, so common knowledge need not be replicated
across routed experts. \Cref{tab:attnalt:moe} lists published routing setups.

\begin{table}[t]
\centering\small
\begin{tabular}{@{}lrrrl@{}}
\toprule
Model & routed experts & active per token & shared & notes\\
\midrule
Switch Transformer & up to 2048 & 1 & 0 & top-1 routing\\
Mixtral 8x7B & 8 & 2 & 0 & 47B total, 13B active\\
OLMoE & 64 & 8 & 0 & fine-grained\\
DeepSeek-V2 & 160 & 6 & 2 & fine-grained + shared\\
DeepSeek-V3 & 256 & 8 & 1 & aux-loss-free balancing, 671B total / 37B active\\
Qwen3-30B-A3B & 128 & 8 & 0 & see \cref{sec:modern:moe}\\
\bottomrule
\end{tabular}
\caption{Expert configurations of some published MoE models (per MoE layer).}
\label{tab:attnalt:moe}
\end{table}

\begin{figure}[t]
\centering
\begin{tikzpicture}[
  ex/.style={draw, rounded corners=1pt, minimum height=0.55cm, font=\tiny, inner sep=1pt},
  arr/.style={-{Stealth[length=1.4mm]}, thick}]
\node[font=\scriptsize\bfseries, align=center] at (1.75,0.95) {coarse:\\8 experts, top-2};
\foreach \i in {0,...,7} {
  \pgfmathsetmacro{\c}{(\i==2 || \i==5) ? 45 : 8}
  \node[ex, minimum width=0.42cm, fill=sysblue!\c] (a\i) at (\i*0.5,0) {};
}
\node[font=\scriptsize] at (1.9,-0.65) {each expert: $d_\text{ff}$ wide};
\begin{scope}[xshift=6.3cm]
\node[font=\scriptsize\bfseries, align=center] at (2.3,0.95) {fine-grained + shared:\\32 routed, top-8, 1 shared};
\node[ex, minimum width=0.42cm, fill=keygreen!45] at (-0.6,0) {S};
\foreach \i in {0,...,31} {
  \pgfmathsetmacro{\c}{(\i==1 || \i==4 || \i==9 || \i==13 || \i==18 || \i==22 || \i==27 || \i==30) ? 45 : 8}
  \node[ex, inner sep=0pt, minimum width=0.11cm, fill=sysblue!\c] at (\i*0.15,0) {};
}
\node[font=\scriptsize] at (2.3,-0.65) {each routed expert: $d_\text{ff}/4$ wide};
\end{scope}
\end{tikzpicture}
\caption{Fine-grained experts with a shared expert (right) activate the same number of MLP FLOPs as a
coarse top-2 design (left) but choose among $\binom{32}{8}\approx10^7$ combinations instead of
$\binom{8}{2}=28$. The always-on shared expert (S) holds common knowledge.}
\label{fig:attnalt:moe}
\end{figure}

\paragraph{Keeping routers stable.} Routers are small softmax classifiers, and their logits are prone to the
same growth as LM logits. ST-MoE (Zoph et al., 2022) computes the router in fp32 and adds a \emph{router
z-loss} on $\log\sum_e\exp(r_e)$, analogous to \cref{eq:archdesign:zloss}. Older recipes add multiplicative
jitter noise to router inputs to make experts less brittle.

\paragraph{Balancing without an auxiliary loss.} The auxiliary loss of \cref{eq:modern:aux} trades
quality for balance. DeepSeek-V3 instead adds a per-expert bias $b_e$ to the routing scores used only for
\emph{selection} (not for the combination weights) and nudges it after each step: decrease $b_e$ if expert
$e$ was overloaded and increase it if underloaded. This is an online control loop that keeps loads balanced
with only a tiny sequence-level auxiliary term.

\paragraph{Capacity, dropping, and nondeterminism.} With a fixed capacity per expert
(\cref{sec:modern:moe}), tokens beyond capacity are dropped and pass through on the residual stream
only. Which tokens are dropped depends on everything else in the batch, so in a serving system \emph{another
user's request can change your output}. Dropless implementations such as MegaBlocks use block-sparse or
grouped GEMMs to process variable-size expert batches without capacity limits.

\begin{systemslens}[Dispatching tokens to experts]
An MoE layer's runtime is dominated by data movement and load imbalance, not by FLOPs. With expert
parallelism (\cref{ch:distributed}) each token's hidden vector is sent to the devices holding its $k$ experts
and the results are sent back: two all-to-all exchanges of $k\cdot d$ values per token per MoE layer. For
DeepSeek-V3-like shapes ($k=8$, $d=7168$, bf16) that is about 115\,KB per token each way, which is why it
restricts each token to experts on at most four nodes and overlaps communication with computation. On each
device, the tokens routed to its experts form groups of different sizes; a \emph{grouped GEMM} (one kernel
launch over many small matrices of varying row counts) or block-sparse formulation, as in MegaBlocks,
processes them without padding to a capacity. The slowest expert sets the step time, so balance is a
throughput problem as much as a modelling one.
\end{systemslens}

\paragraph{Upcycling and fine-tuning.} An MoE can be initialised from a dense checkpoint by copying its MLP
into every expert and adding a fresh router (``sparse upcycling'', Komatsuzaki et al., 2023). Several
released MoEs were reportedly built this way. MoEs fine-tuned on small datasets can overfit, which some
recipes mitigate by fine-tuning only non-expert parameters or by using much larger SFT sets.

% =====================================================================
\section{Choosing a sequence mixer}\label{sec:attnalt:choose}

\begin{table}[t]
\centering\small
\begin{tabularx}{\linewidth}{@{}l>{\raggedright\arraybackslash}X>{\raggedright\arraybackslash}X>{\raggedright\arraybackslash}X>{\raggedright\arraybackslash}X@{}}
\toprule
Mixer & Training cost & Decode state & Exact long-range recall & Ecosystem\\
\midrule
Softmax (GQA) & $O(T^2)$ FLOPs, IO-optimal kernels & grows with $T$ & strong & mature everywhere\\
Sliding window & $O(Tw)$ & bounded by $w$ & only within window & mature\\
MLA & $O(T^2)$ & grows, $\sim$57$\times$ smaller & strong & engine support growing\\
Linear / GLA / Mamba-2 & $O(TCd)$, chunked matmuls & constant & weak (capacity-limited) & specialised kernels\\
Gated DeltaNet & chunked, WY form & constant & better than additive & specialised kernels\\
Hybrid (1:3--1:7) & mostly linear & small + constant & close to full attention & needs dual cache management\\
\bottomrule
\end{tabularx}
\caption{A qualitative comparison of sequence mixers. ``Exact recall'' refers to tasks such as associative
recall and needle retrieval; the hybrid entry summarises reported results, not a controlled comparison.}
\label{tab:attnalt:choose}
\end{table}

\Cref{tab:attnalt:choose} summarises the trade-offs. For a new model whose deployment is dominated by long
generations, a hybrid with a small fraction of full-attention layers is now a credible default. For a model
that must be served everywhere tomorrow, plain GQA attention with sliding-window layers remains the
lowest-risk choice, because every kernel library and inference engine supports it.

% =====================================================================
\section{Summary}\label{sec:attnalt:summary}
\begin{itemize}
\item Long-context cost is dominated by the KV cache at decode time; shrinking it (GQA, MLA, quantisation),
  bounding it (windows, sinks) or replacing it (recurrent states) are the three strategies.
\item Removing the softmax makes attention associative: the same operator has parallel ($O(T^2)$),
  recurrent ($O(T)$, constant state) and chunkwise forms. Training uses the chunked form and decoding the
  recurrent one.
\item Decay ($\gamma$) adds forgetting: constant per head (RetNet) or data-dependent (GLA, Mamba-2). It turns
  the causal mask into a decay mask.
\item The delta rule writes errors instead of sums, giving overwrite semantics and a test-time-training
  interpretation; gated DeltaNet adds forgetting.
\item Selective SSMs are first-order linear recurrences, parallelisable with associative scans, and run
  efficiently only with IO-aware kernels or chunked-matmul reformulations.
\item Hybrids keep a few attention layers for exact retrieval; they complicate caching, prefix reuse and
  speculative decoding in serving engines.
\item Modern MoEs use fine-grained and shared experts, fp32 routers with z-loss, bias-based balancing and
  dropless kernels; capacity-based dropping makes outputs batch-dependent.
\end{itemize}

% =====================================================================
\section{Exercises}\label{sec:attnalt:exercises}

\begin{exercise}[Normaliser for free]
Show that appending a column of ones to $\mat{V}$ in any of the three linear-attention forms yields the
denominator of \cref{eq:attnalt:linear} in the last output column. Why can this denominator approach zero
with an unbounded feature map, and why does $\mathrm{elu}+1$ avoid sign problems but not small values?
\end{exercise}

\begin{exercise}[Chunked derivation]
Starting from \cref{eq:attnalt:gated}, derive the chunked form with decay: express $\vect{y}_i$ for $i$ in chunk
$c$ in terms of the state at the end of chunk $c-1$ and the tokens of chunk $c$, and give the state update.
Verify that your expressions match \cref{lst:attnalt:chunked}.
\end{exercise}

\begin{exercise}[Crossover]
For a model with $d=4096$, 32 heads of $\dhead=128$ and chunk size 64, plot per-token forward FLOPs of causal
softmax attention and chunked linear attention against $T$. At $T=32{,}768$, what is the ratio?
\end{exercise}

\begin{exercise}[Capacity of a matrix memory]
Write $n$ random unit keys in $\R^{d}$ with additive updates $\mat{S}=\sum_i\vect{k}_i\vect{v}_i^\top$ and
retrieve with $\vect{q}=\vect{k}_j$. Derive the signal-to-interference ratio as a function of $n/d$. Measure
it with \code{llmsys.linear\_attn}, then repeat with the delta rule ($\beta=1$, sequential writes) and
explain the difference.
\end{exercise}

\begin{exercise}[Scan work and depth]
Count the combine operations and parallel rounds used by Hillis--Steele and by Blelloch scans for $T=2^{16}$.
Which would you choose on a GPU with 132 SMs for a single sequence? For a batch of 256 sequences?
\end{exercise}

\begin{exercise}[Underflow]
A gate of $\gamma=0.9$ applied over 1000 steps gives $0.9^{1000}\approx10^{-46}$, which underflows bf16 and
fp16. Identify the operations in \cref{lst:attnalt:chunked} that would fail if implemented with explicit
products and divisions, and explain why working with differences of cumulative log-gates is safe.
\end{exercise}

\begin{exercise}[Hybrid budget]
For the Qwen3-8B shape of \cref{fig:attnalt:cache}, find the largest attention fraction (one full layer every
$k$ layers) such that a batch of 32 sequences of 128k tokens fits in 40\,GB of cache memory alongside a
sliding-window alternative for the remaining layers with $w=4096$.
\end{exercise}

\begin{exercise}[MLA absorption]
Show algebraically that $\vect{q}^\top\mat{W}^{UK}\vect{c}$ can be computed as $(\mat{W}^{UK\top}\vect{q})^\top\vect{c}$
and count the FLOPs of both orders during decode with $H$ heads, $\dhead=128$ and $d_c=512$. Why does
absorption help decode but not necessarily training?
\end{exercise}

\begin{exercise}[Batch-dependent MoE outputs]
Construct a small example of top-1 routing with capacity factor 1.0 where the output for token A changes when
token B is added to the batch. What does this imply for testing an MoE inference engine for determinism?
\end{exercise}

\begin{exercise}[Serving a recurrent layer]
Design the data structures a continuous-batching engine (\cref{ch:serving}) needs for a hybrid with paged KV
for attention layers and fixed-size state slots for recurrent layers. How do you support (a) prefix caching
of a shared system prompt, (b) parallel sampling of 8 continuations, and (c) preemption with swap-out?
\end{exercise}

% =====================================================================
\section{Lab 8: Linear attention, the delta rule, and selective scans}\label{sec:attnalt:lab}
\begin{labbox}{8}{Linear attention, the delta rule, and selective scans}
\labdir{lab08_attnalt}
\labsection{Goals}
Build the core operators of non-softmax sequence mixers in their parallel, recurrent and chunked forms;
prove their equivalence with tests; and quantify their memory advantage at decode time.

\labsection{Setup}
\begin{lstlisting}[style=shell]
cd code
LAB_IMPL=starter  python -m pytest -q labs/lab08_attnalt
LAB_IMPL=solution python -m pytest -q labs/lab08_attnalt
python labs/lab08_attnalt/bench.py --device cpu     # or --device cuda
\end{lstlisting}
Reference layers: \code{llmsys/linear\_attn.py} (\code{LinearAttention}, all three forms, delta rule,
memory calculators) and \code{llmsys/ssm.py} (\code{scan\_parallel}, \code{MambaBlock}). Tests use float64
so equivalences can be checked to $10^{-10}$.

\labsection{Tasks}
\begin{enumerate}
\item \code{linear\_attention\_quadratic}: masked $(\mat{Q}\mat{K}^\top\odot\mat{D})\mat{V}$ with a decay
  mask built from cumulative log-gates; support no decay, per-head decay $(H,)$ and per-step gates $(B,H,T)$.
\item \code{linear\_attention\_recurrent}: the RNN form, returning the final state so that a sequence can be
  continued later (tested by splitting a sequence in two).
\item \code{linear\_attention\_chunked}: the chunkwise form; it must agree for chunk sizes 1, 4, 7 and 64.
\item \code{delta\_rule\_step}: the (gated) delta update. The test checks exact overwrite of repeated keys.
\item \code{associative\_scan}: Hillis--Steele scan for $h_t=a_th_{t-1}+b_t$.
\item \code{selective\_ssm\_step}: one discretised Mamba step; the test plugs it into \code{MambaBlock} and
  compares against the block's own forward pass.
\item \code{sliding\_window\_mask} and \code{decode\_cache\_bytes} for full, windowed, linear and hybrid stacks.
\end{enumerate}

\labsection{Measurements}
\begin{itemize}
\item Run \code{bench.py} at increasing context lengths. On a GPU, per-token decode time of softmax
  attention should grow linearly with $T$ while the linear and SSM steps stay flat; find the crossover length.
  (On a busy CPU the numbers are noisy; look for the trend over a wide range of $T$.)
\item Reproduce \cref{fig:attnalt:cache} with \code{decode\_cache\_bytes} for a model of your choice.
\item Train a two-layer model on the copy lines of \code{llmsys.ablate.make\_corpus} with (a) softmax
  attention, (b) gated linear attention, (c) a 1:1 hybrid. Compare the loss on the copied segment only.
\end{itemize}

\labsection{Deliverables}
Passing tests; a decode-time plot; a cache-memory plot; a paragraph explaining the copy-task results in terms
of state capacity.

\labsection{Stretch goals}
Implement the Blelloch (work-efficient) scan; implement a chunked delta rule using the WY representation and
test it against the recurrent form; write a Triton kernel for the chunked forward that keeps $\mat{S}$ in
shared memory (\cref{app:triton}).
\end{labbox}
